% ------------------------------------------------------------
% Page 82
% ------------------------------------------------------------

traqué dépêche toi de lui régler son compte ». ce que je fis bien que l’on fût un dimanche.
Effectivement 2 jours après nous eûmes la visite des gendarmes qui recherchaient ce garçon.

\vspace{0.45cm}

L’usine commençait à tourner, mais dans des conditions épouvantables. Les pannes se
succédaient. Il fallait toute l’astuce et l’imagination de Rolland et de son adjoint pour arriver à
résoudre les problèmes les plus urgents. Il fallait descendre à Nîmes avec un mouton pour
avoir des pneus synthétiques, un gigot pour avoir des pointes ou des pièces de rechange. Les
aciers ne valaient rien ; tout cassait. La route allant en forêt était épouvantable. Les chauffeurs
dont notre ami Persico descendaient leurs 5 tonnes de bois à raison de 2 voyages par jour dans
des conditions inimaginables.

\vspace{0.45cm}

Au début de l’automne 1943 l’Hôtel de France à Meyrueis fut réquisitionné par la Préfecture.
On vit arriver un groupe de polonais placés en résidence surveillée. Le STO nous autorisa à
embaucher ceux qui voulaient travailler. Nous ignorions les origines de ces gens, 7 ou 8
vinrent travailler. Je me souviens de l’un d’eux d’assez forte corpulence qui s’appelait Pakula.
Il parlait assez bien le français.

\vspace{0.45cm}

Vers le milieu de l’hiver 1944, M. Rolland qui revenait d’un voyage à Nîmes me dit qu’il
venait d’embaucher un jeune technicien, élève à l’école des Arts et Métiers qui s’appelait
« Bonnet » ainsi qu’un autre qui répondait au nom de « Séguret ». Une heure après alors que
j’étais seul dans mon bureau Rolland vint me dire à mi-voix et en me regardant bien dans les
yeux : Caban el ! Bonnet c’est mon frère prisonnier évadé d’Allemagne et Séguret c’est mon
cousin germain. Ce dernier requis par le STO s’était accidentellement cassé une jambe en
Allemagne. Après une hospitalisation de 2 mois il avait bénéficié d’un congé de
convalescence de 10 jours. Sur le point de repartir Rolland l’avait engagé à venir se camoufler
à l’usine de Meyrueis. Bien sûr, j’assurai Rolland de la plus grande discrétion. Peu après
j’assistai à la rencontre des deux cousins dans mon bureau, l’un demandant à l’autre
« Comment t’appelles-tu maintenant ? ».

\vspace{0.45cm}

Des parents du PDG de la Cévenole qui avaient une vingtaine d’années vinrent se faire
embaucher pour éviter le STO. Ils arrivaient de Nîmes qui était occupée par les Allemands.
D’autres, du même âge, venaient d’Avignon. Il y avait le fils d’un armurier de cette ville.

\vspace{0.55cm}

\textbf{Réfractaires et Etrangers à l’usine}

\vspace{0.35cm}

D’autre part il y avait un jeune mécanicien appelé Ausset. Son père était chef de service au
Ravitaillement à Nîmes. Par lui Rolland put se procurer des cartes d’alimentation pour les
irréguliers. Il y avait aussi un contre maître nommé Carrier Victor venu de Chapareillan pour
éviter le STO.

\vspace{0.45cm}

Plus tard le PDG de Nîmes nous envoya un autre réfractaire, Joseph Mondino, d’origine
niçoise. Il travaillait comme metteur au point des moteurs d’avion à Marignane. Les
allemands ayant envahi son usine, il put miraculeusement s’échapper et fut dirigé sur
Meyrueis. Nous avons à ce jour conservé d’excellentes relations d’amitié. Joseph Mondino
déployait des trésors d’imagination pour faire tourner les moteurs fatigués. Le contremaître
forestier s’appelait Neufeld. Il était marié, originaire d’Obernai en Alsace. Il s’était échappé
et avait pu, après avoir effectué un temps aux Chantiers de Jeunesse, se faire embaucher.

\vspace{0.45cm}

Rolland qui a tout fait pour rendre service ,en camouflant tous ceux qui étaient traqués, en
était arrivé à la situation suivante :
